% Complete authored Preface, original PDF 9-10 / printed ix-x.
% block: preface-heading page:9 type:section
\section*{{前言}}\phantomsection\label{sec:preface}

% block: preface-computational-methods page:9 type:ordinary-paragraph
{求微分方程近似解的计算方法在科学、工程、数学和技术中起着至关重要的作用. 事实上, 支配物理世界的关键过程, 包括波的传播、热力学、流体流动、固体变形、气体动力学、电磁现象、量子力学、广义相对论以及许多其他过程, 都由微分方程描述; 我们依靠数值方法来模拟、探索、预测和控制涉及这些过程的系统. 这些应用中出现的微分方程问题种类繁多, 人们为发展能够准确、高效地求解不同问题的数值方法进行了大量研究. 对这些算法的数学分析起着不可或缺的作用: 在明确界定的情形下为特定方法提供严格论证, 给出定量误差界, 并使不同方法之间的比较成为可能. 本书着重讨论有限元方法. 这是一大类求解微分方程的数值方法, 适用范围广、用途大; 与此相应, 关于这类方法也有大量数学分析成果, 这并非偶然.}

% block: preface-functional-analysis-and-structure page:9 type:ordinary-paragraph
{在计算科学的许多领域, 有限元方法在实践和理论上都是成熟的工具. 尽管如此, 所出现的偏微分方程 (PDE) 种类繁多, 仍有许多重要问题使现有数值途径失效, 有效的数值方法尚有待构造. 因此, 致力于构造和分析新方法的研究蓬勃发展. 传统上, 数值 PDE, 特别是有限元方法研究所用的关键数学工具来自泛函分析: Hilbert 空间和 Banach 空间、Hilbert 投影定理、Lax–Milgram 引理、Bramble–Hilbert 引理、对偶性、Sobolev 空间等. 本书介绍的有限元外微积分 (FEEC) 同样实质性地依赖泛函分析, 尤其是 Hilbert 空间上闭无界算子的理论. 但是, FEEC 的数学工具远不止于泛函分析: 它引入几何和拓扑的工具, 为那些较传统途径难以处理的 PDE 类发展和分析数值方法. 从 FEEC 导出的方法是保结构数值方法的典型例子, 因为这些方法旨在离散层次保留 PDE 的关键几何、拓扑和代数结构. 事实表明, 对于标准方法失效的多种问题, 这是发展有效有限元方法的关键. 具体而言, FEEC 着眼于与 Hilbert 函数空间上微分算子复形相关的 PDE, 使用构成这些复形的子复形、并且能通过交换投影与原复形联系起来的有限元空间.}

% block: preface-history page:10 type:ordinary-paragraph
{虽然 FEEC 的先驱工作可以追溯到几十年前, 即有限元方法的早期甚至更早, 但它作为一种独立理论, 最初开始形成于我在 2002 年国际数学家大会上的报告 \BookCite{5}, 并在我与 Richard Falk 和 Ragnar Winther 合作撰写的两篇长文, 即 2006 年的 \BookCite{11} 和 2010 年的 \BookCite{13}, 中得到了形式化表述. 第一篇论文着重研究一个特定的微分算子复形, 即 $\mathbb R^3$ 中区域 (或 Riemannian 流形) 上微分形式的 de Rham 复形. \emph{有限元外微积分}这一名称正是在那里首次出现的, 其中的外微积分指微分形式的微积分. 在 2010 年的论文中, 我们更加重视 Hilbert 复形的抽象结构; $L^2$ de Rham 复形是它的一个特例, 这使 FEEC 能够处理其他应用中出现的其他复形.}

% block: preface-course-and-audience page:10 type:ordinary-paragraph
{到 2012 年 6 月, FEEC 理论的基本轮廓已经形成. 我有幸获得机会, 为来自世界各地的近 70 名教师、研究生和其他研究人员开设了一门 FEEC 集中短期课程. 作为 NSF-CBMS 系列会议的一部分, 该课程得到国家科学基金会和数学科学会议委员会的慷慨资助, 并由 Brown University 的数学计算与实验研究所 (ICERM) 出色地承办. 本书由这门课程发展而来. 它与课程有着相同的目标: 帮助数值分析工作者在一处快速掌握 FEEC 的基础, 包括所需的几何和泛函分析预备知识. 但本书面向的读者比课程更广, 还希望数值分析之外其他领域的数学家和数学专业学生也能阅读, 只要他们有兴趣了解几何和拓扑技术如何在数值 PDE 中发挥作用. 课程之后, FEEC 一直在蓬勃发展, 因此本书包含的内容远多于课程所讲的内容, 其中一些内容在当时甚至尚未形成.}

% block: preface-book-organization page:10 type:ordinary-paragraph
{本书的第一部分, 即 Chapters~\ref{chap:introduction}--\ref{chap:approximation-hilbert-complexes}, 简要而迅速地建立所需的同调代数、代数拓扑和泛函分析预备知识. 这些要素在 Chapter~\ref{chap:hilbert-complexes} 所研究的 Hilbert 复形这一基本结构中结合起来. 仍在这一一般抽象框架内, Chapter~\ref{chap:approximation-hilbert-complexes} 建立与 Hilbert 复形相关的问题的逼近理论. 本书第二部分由 Chapters~\ref{chap:exterior-calculus} 和~\ref{chap:finite-element-differential-forms} 组成, 将一般理论应用于 Hilbert 复形最典型的例子, 即 $\mathbb R^n$ 中区域上的 $L^2$ de Rham 复形. 最后, 在结尾一章中, 我们简要介绍其他一些例子和应用.}

% block: preface-acknowledgments page:10 type:ordinary-paragraph
{感谢 NSF 和 CBMS 对 2012 年 FEEC 课程及本书的支持, 也感谢在我发展 FEEC 和撰写本书期间 NSF 项目 DMS-1115291、DMS-1418805 和 DMS-1719694 所给予的资助. 对 CBMS 的前任和现任主任 Ron Rosier 和 David Bressoud 的耐心和理解, 尤应致谢. 还要感谢 ICERM 承办这门课程, 尤其感谢构思并组织课程的 Alan Demlow、Johnny Guzmán 和 Dmitriy Leykekhman. 课程的听众同样给了我很大的激励和启发, 其中许多人后来对 FEEC 作出了重要贡献. 还有一些人通读了全部或部分书稿, 提出了无数宝贵建议: 感谢 Richard Falk、Ragnar Winther、Shawn Walker、Espen Sande 和 Kaibo Hu. Johnny Guzmán、Anil Hirani 和 Ragnar Winther 甚至把本书的早期版本用作他们所授课程的教材, 这尤其有帮助.}
